% Transcription — fonds Alexandre Grothendieck, shelfmark 41, batch 1 (pages 1-13).
% Established from the facsimile archives/batches/41/batch-01.pdf (a mirror of
% https://grothendieck.umontpellier.fr/41.pdf), per the skill
% transcribe-grothendieck. The folder is thirteen pages and this is its only
% batch.
%
% Twelve of the thirteen leaves are transcribed. Page 1 is the title leaf: it
% carries, alone at the head of an otherwise blank sheet, an underlined
% « Chap. VI », one line of prose, two struck words, and the run's title
% underlined in full. Because it carries more than a name it keeps its \page
% marker, and the title is also lifted out as the first section heading. Page 2
% is the reverse of that leaf and carries nothing; it is skipped, and the gap
% in the page numbers is the record. Pages 3 to 13 are the run.
%
% The subject is the one the title names: an étale Y-scheme M indexing the
% relatively irreducible components of the fibres of a morphism f : X -> Y, and
% the notion of a « monogène » section that goes with it. Page 3 defines
% monogène for X/Y and for a map phi : Y -> X, forms the functor F(Y') of
% monogène sections after base change, asks whether it is representable, and
% gives the criterion an (M, phi_M) must satisfy to represent it: over an
% algebraically closed extension Omega of kappa(y), M(Omega) must be in
% bijection with the set of irreducible components of X tensor_Y Omega. Page 4
% reduces the existence to a local question — the problem is local on Y, one
% may take Y = Spec(A) with A local, then complete by descent of étale
% morphisms. Pages 5 and 6 run the construction under the hypotheses
% « universellement ouvert, de type fini, A local noethérien complet,
% f^{-1}(y) irréductible »: a point x of the closed fibre determines a unique
% relatively irreducible component, that map phi is monogène, and the last
% thing to check is that phi is continuous. Page 7 states the Lemme the
% construction rests on, with its five hypotheses, and closes on the
% separatedness of M/Y. Page 8 is the Théorème: under a local
% quasi-finiteness-over-affine-space hypothesis — satisfied in particular when
% X is flat over Y — an étale separated M indexing the relatively irreducible
% components of X/Y exists if and only if two specialisation conditions hold.
% Page 9 adds the N.B. that the components of the geometric fibre inject into,
% but need not exhaust, those of the fibre, and a third condition (iii) for M
% to be an open étale scheme. Page 10 is two Remarques — (i) alone gives a
% non-separated M; (iii) is automatic for a separable f; and if f is proper, M
% is proper. Page 11 backs up to the irreducible-fibre case: Prop. 1 lists
% three equivalent conditions for phi to be continuous, adds the auxiliary
% (iii bis), and then Déf. 2 and Déf. 3 fix « universellement monogène » and
% « monogène » once and for all. Page 12 observes that all of it is local on Y,
% asks whether the functor I(X xS S'/S') is representable — « Est-il
% représentable ??? » — and draws the cube of base changes in pencil. Page 13
% is a Lemme cancelled by two long diagonals, two base-change diagrams, and a
% Hyp/Concl pair on specialisations of relatively irreducible components, the
% second conclusion Concl' being struck out in turn.
%
% The hand is drafting throughout and that is the folder's governing fact. The
% statements, the labelled conditions, the displayed formulas and the
% underlined words carry; a large part of the connective prose does not, and it
% is marked \ill{} rather than invented. The apparatus is dense in consequence
% and that density is the record of what these leaves support.
%
% Notation fixed once here rather than page by page. He writes two adverbs as
% an abbreviation closed by a raised, underlined « t »: « rel^t » for
% « relativement » and « abs^t » for « absolument », both set here as
% $\mathrm{rel}^{\mathrm{t}}$ and $\mathrm{abs}^{\mathrm{t}}$; the pair is
% settled by page 9, where he strikes a word and writes « relatives » above it,
% and by page 8, where « abs^t » is inserted interlinearly over
% « irréductibles ». E_{X/Y}(y) is his notation for the set of relatively
% irreducible components of the fibre over y, and it is used unexplained from
% page 5 on. « f. et s. » is « il faut et il suffit ».
%
% Readings flagged and not settled: the two words after « Dans le § de la » on
% page 1; the noun in « une Y-préschéma » on page 7, whose ending will not
% resolve; the underlined heading above « Supposons » on the same page, read
% « Le Moyen »; and the underlined word opening the refinement of page 9, read
% « Cor. ». One divergence from the inventory is recorded where it occurs and
% not repaired: the leaf reads « Chap. VI », the archivists' title « (Chap. V) ».
%
% Pass: Opus 5 (claude-opus-5), 2026-09-13 — first pass, unchecked against the pages by a human.
\documentclass[11pt,a4paper]{article}
% Le préambule commun à toutes les transcriptions du fonds Grothendieck.
%
% Il tient en une idée : **l'incertitude se compose, elle ne se résout pas.**
% Une transcription qui lisse un mot douteux a détruit la seule chose pour
% laquelle on la faisait — un lecteur doit pouvoir savoir, à chaque endroit,
% ce qui était sur le feuillet et ce qui a été ajouté. Les six macros qui
% suivent sont donc l'appareil critique, et non de la décoration.
%
% Les mêmes commandes sont reconnues par `scripts/render.mjs`, qui produit la
% vue de lecture du site. Une macro ajoutée ici doit l'être aussi là-bas,
% faute de quoi le rendu échouera bruyamment — ce qui est voulu : un rendu
% silencieusement faux est pire qu'un rendu absent.

\ProvidesPackage{grothendieck}[2026/08/08 Transcriptions du fonds Grothendieck]

\usepackage{amsmath,amssymb,amsthm}
\usepackage{mathrsfs}
\usepackage{stmaryrd}
% Les diagrammes commutatifs sont partout dans ce fonds : c'est la langue
% même de la géométrie algébrique de Grothendieck, et une transcription qui
% les rendrait en prose ne transcrirait rien.
\usepackage{tikz-cd}
% Français par défaut — c'est la langue des feuillets — mais l'anglais est
% chargé aussi : la traduction et le résumé sont des documents anglais, et la
% typographie française (espaces avant les deux-points) y serait une faute.
% Ces documents basculent par \selectlanguage{english} en tête de corps.
\usepackage[english,french]{babel}
\usepackage{fontspec}
\usepackage{geometry}
\geometry{a4paper,margin=25mm,marginparwidth=28mm}
\usepackage{marginnote}
\usepackage{xcolor}
% ulem, not soul. `soul`'s \st and \ul are built on a character-by-character
% scan that XeLaTeX's font machinery breaks: the document fails with an
% undefined control sequence rather than degrading. ulem does the same two jobs
% and is Unicode-engine safe. `normalem` stops it hijacking \emph, which the
% transcriptions use for Grothendieck's own emphasis.
\usepackage[normalem]{ulem}
\usepackage{hyperref}
\hypersetup{colorlinks=true,linkcolor=black,urlcolor=black,pdfborder={0 0 0}}

% --- Métadonnées du lot -----------------------------------------------------
% Reprises telles quelles par le script de rendu, qui les affiche en tête de
% la vue de lecture. Elles fixent ce que le fichier prétend couvrir : sans
% elles, une transcription flotte sans point d'attache dans un fonds de seize
% mille feuillets.

\newcommand{\folder}[1]{\gdef\grothFolder{#1}}
\newcommand{\batch}[1]{\gdef\grothBatch{#1}}
\newcommand{\leaves}[2]{\gdef\grothFirst{#1}\gdef\grothLast{#2}}
\newcommand{\foldertitle}[1]{\gdef\grothFoldertitle{#1}}
\gdef\grothFolder{?}\gdef\grothBatch{?}\gdef\grothFirst{?}\gdef\grothLast{?}\gdef\grothFoldertitle{}

% --- Le résumé --------------------------------------------------------------
%
% Il ouvre la lecture modernisée, sous le titre « Résumé », et il est composé
% en petit. Ce n'est pas un ornement : c'est le seul passage du document qui
% ne soit pas des mathématiques, et un lecteur doit pouvoir le reconnaître
% avant de l'avoir lu — puis le sauter s'il connaît déjà le sujet, ou s'y
% arrêter s'il ne le connaît pas. Le filet qui le suit marque la reprise du
% corps.

\newenvironment{resume}
  {\small\setlength{\parskip}{0.45em}}
  {\par\medskip\hrule\medskip}

% --- Mots-clés ---------------------------------------------------------------
%
% En anglais, en clôture du résumé de la lecture modernisée : le vocabulaire
% moderne sous lequel chercher ce que les feuillets construisent. Le manifeste
% du site (`npm run manifest`) les extrait pour étiqueter la cote entière —
% cette ligne est la source unique des tags, il n'y a pas de fichier à côté.
\newcommand{\keywords}[1]{\par\smallskip\noindent
  {\footnotesize\textsc{Keywords} --- \textit{#1}}\par}

% --- L'appareil critique ----------------------------------------------------

% Le numéro de feuillet, en marge. C'est l'ancre qui permet de régler un
% désaccord sur une formule en regardant *un* feuillet plutôt qu'une chemise
% de six cents. Le site s'en sert aussi pour tourner les pages du fac-similé
% quand on fait défiler la transcription.
\newcommand{\leaf}[1]{%
  \marginnote{\footnotesize\color{gray!70}\textsf{#1}}%
  \label{leaf:#1}\ignorespaces}

% Illisible : on ne devine pas. Un mot inventé qui se lit comme les autres est
% le pire résultat possible.
\newcommand{\ill}{\textcolor{red!60!black}{[\,\ldots\,]}}

% Lecture douteuse : le mot est donné, et signalé comme incertain.
\newcommand{\uncertain}[1]{\uline{#1}}

% Ajout de l'éditeur — une lettre manquante, un mot sous-entendu.
\newcommand{\add}[1]{\textcolor{blue!50!black}{[#1]}}

% Biffé par Grothendieck lui-même. Ce qu'il a rayé fait partie du document :
% une rature dit souvent où la pensée a changé de direction.
\newcommand{\struck}[1]{\sout{#1}}

% Note du transcripteur, sur l'état matériel du feuillet.
\newcommand{\note}[1]{\par\smallskip{\small\color{gray!60!black}\textit{[#1]}}\par\smallskip}

% Note marginale de Grothendieck, distincte du corps du texte.
\newcommand{\marginal}[1]{\par\smallskip\noindent\hspace{1em}%
  {\small\color{gray!60!black}#1}\par\smallskip}

% --- Titre ------------------------------------------------------------------

\AtBeginDocument{%
  \begin{center}
    {\large\bfseries Fonds Alexandre Grothendieck --- cote n\textsuperscript{o}~\grothFolder}\\[2pt]
    {\itshape\grothFoldertitle}\\[4pt]
    {\small feuillets \grothFirst--\grothLast\ (lot \grothBatch)}\\[2pt]
    {\footnotesize\color{gray!60!black}Transcription établie d'après le fac-similé numérisé
     par l'Université de Montpellier.\\
     Les lectures incertaines sont soulignées, les illisibles notées [\,\ldots\,],
     les ajouts entre crochets.}
  \end{center}
  \vspace{1em}}

\endinput


\folder{41}
\batch{1}
\pages{1}{13}
\dating{[années 1960-1970]}
\watermark{Édition de démonstration}
\foldertitle{Construction de morphismes étales (Chap. V) : notes manuscrites (s.d.)}

\begin{document}

\section{Morphismes étales associés à l'étude des composantes irréductibles d'un morphisme ouvert}

\note{le titre est de lui : il est écrit, souligné, sur le feuillet 1, par
ailleurs vide, sous un « Chap. VI » souligné et sous deux mots biffés dont le
premier, « Construction », est celui que les archivistes ont retenu pour leur
propre titre}

\page{1}

Chap. VI. Dans le § de la \uncertain{descente}, \ill{}, applications,

\struck{Construction} \struck{Étude}

\note{suivent, soulignés et occupant deux lignes, les mots qui donnent son
titre à la section : « Morphismes étales associés à l'étude des composantes
irréductibles d'un morphisme ouvert ». Le reste du feuillet est vide}

\note{le chiffre du chapitre est un VI sur le feuillet, et non le V que porte
l'inventaire ; le groupe où les archivistes rangent la cote est celui du
Séminaire de Géométrie Algébrique}

\page{3}

$X$ de type fini sur $Y$. On dit que $X$ est \emph{monogène} sur $Y$ si

\begin{enumerate}
\item[(i)] les fibres de $X$ sur les $y \in Y$ sont
$\mathrm{rel}^{\mathrm{t}}$ irréductibles ;
\item[(ii)] $X \to Y$ est ouvert.
\end{enumerate}

\note{la condition (ii) est ajoutée en interligne au-dessus de la ligne de (i),
d'une écriture plus fine}

Une application $\varphi : Y \to X$ est dite \emph{monogène} s'il existe un
ouvert $U$ de $X$, \add{universellement ouvert} sur $Y$, tel que pour tout
$y \in Y$, $\varphi(y)$ soit le point générique de $U \cap f^{-1}(y)$.

\note{le mot entre « un ouvert $U$ de $X$ » et « tel que » est repris deux
fois ; la seconde main donne « universellement ouvert sur $Y$ », qui est aussi
ce que la définition 3 de la page 11 retiendra}

\add{Si $Y$ est \ill{} noethérien, il revient au même de dire :}

\begin{enumerate}
\item[(ii)] Pour tout $y \in Y$, $\varphi(y)$ est le point générique d'une
composante $\mathrm{rel}^{\mathrm{t}}$ irréductible de $f^{-1}(y)$ ;
\item[(iii)] \ill{}
\end{enumerate}

\marginal{\ill{} a') $\Rightarrow$ (ii) \ill{}}

\note{bloc marginal en biais le long du bord gauche, à quelque 45°,
sur cinq lignes ; il n'est pas lisible à la résolution du film au-delà de
l'implication qui l'ouvre}

Si $\varphi : \struck{X} \to \struck{X}$ est monogène, alors pour tout
changement de base $Y' \to Y$, l'application associée $\varphi' : Y' \to X'$
\add{définie de façon évidente} est monogène. Ainsi l'ensemble $F(Y')$ des
sections monogènes de $X'/Y'$ est un foncteur contravariant en $Y'$.

Est-il représentable, \ill{} ? Si oui, le \ill{} $M$ sur $Y$ est
vraisemblablement étale de type fini sur $Y$ \add{mais pas néc. séparé sur $Y$}.

Soit $M$ étale de type fini sur $Y$, et $\varphi_M$ une \struck{section}
application monogène pour $X \times_Y M / M$, qui ait les propriétés suivantes :
pour tout $y \in Y$ et une extension alg. close $\Omega$ de $\kappa(y)$,
l'application
\[
  M(\Omega) \longrightarrow \{\text{composantes irréductibles de } X \otimes_Y \Omega\}
\]
est bijective. Alors $(M, \varphi_M)$ représente le foncteur $F$.

\marginal{\ill{} $M$ correspond aux points génériques des fibres de $X \to Y$}

\page{4}

\add{On est ramené à prouver ceci :} Soit $\psi : Y \to M$ une application telle
que

\begin{enumerate}
\item[(i)] pour tout $y \in Y$, $\kappa(\psi(y)) = \kappa(y)$, i.e.\ $\psi(y)$
est rationnel sur $\kappa(y)$ ;
\item[(ii)] l'application composée $\varphi : Y' \to X'$, \struck{\ill{}}
faisant correspondre à chaque $y \in Y'$ la composante
$\mathrm{rel}^{\mathrm{t}}$ irréd.\ de $f^{-1}(y)$ correspondant \struck{à
$\kappa(y)$} à $\psi(y)$, est monogène.
\end{enumerate}

Alors \struck{\ill{}} $\psi(Y)$ est \struck{\ill{}} ouvert, \struck{\ill{}}
\struck{$\psi(Y) \to Y$} et donc $\psi$ provient d'une section de $M$ sur $Y$ ;
sous cette forme, c'est évident, grâce à la condition de rationalité inscrite
dans la définition de « monogène ».

\textbf{Conséquence.} Si le Pb.\ des modules a une solution pour
$X \times_Y \mathrm{Spec}(\Omega_y)$, il existe un ouvert $U \ni y$ de $Y$ tel
que le Pb.\ des modules ait une solution pour $X|U / U$.

\note{au-dessus de $\mathrm{Spec}(\Omega_y)$, en interligne : « une
$\mathrm{Spec}(\Omega_y)$ » ; la formule est reprise deux fois sur la ligne}

Cela nous ramène, pour l'existence, dans le cas où \struck{$X$}
$Y = \mathrm{Spec}(A)$, $A$ local. \add{On peut donc supposer \ill{} $Y = $
\ill{}} D'ailleurs, puisque\,: la descente des morphismes étales, on est ramené
au cas où $A$ est complet, et on peut supposer que la fibre \ill{}
$f^{-1}(y_0)$ \ill{} se décompose en composantes $\mathrm{rel}^{\mathrm{t}}$
irréductibles.

\page{5}

\begin{tikzcd}[column sep=small, row sep=small]
X \arrow[d, "f"] \\
Y = \mathrm{Spec}(A)
\end{tikzcd}

\add{strictement} universellement ouvert, de type fini \note{« N.B. » entouré
au-dessus de « de type fini »}, $A$ local noethérien complet, $f^{-1}(y)$
(irréductible).

\marginal{« \uncertain{strictement} » : \ill{}}

\note{bloc en biais dans la marge de gauche, sur cinq lignes, glosant
l'adjectif entre crochets de la première ligne ; illisible au-delà du mot
répété}

\begin{enumerate}
\item[(i)] Supposons que \struck{$f^{-1}(y)$ soit \ill{}} \ill{} $\xi$ qui
\ill{} \ill{} de $X$, il y a deux \ill{} dans $E_{X/Y}(y)$ qui \ill{} \ill{}
a) \ill{}
\end{enumerate}

\note{$E_{X/Y}(y)$ désigne, ici et jusqu'à la fin du dossier, l'ensemble des
composantes $\mathrm{rel}^{\mathrm{t}}$ irréductibles de la fibre au-dessus de
$y$ ; la notation n'est nulle part définie}

\ill{} $Y' = \mathrm{Spec}(W)$, $W$ anneau de valuation discrète \add{= une
extension \ill{}}, \struck{\ill{}} \add{$f^{-1}(D')$
$\mathrm{rel}^{\mathrm{t}}$ irréductible}, donc qu'il y a coïncidence des
\ill{} des fibres géométriques en $f^{-1}(y)$. Supposons que ceci soit exclu.

\begin{enumerate}
\item[(ii)] Soit $\xi \in E_{X/Y}(y)$ tel que $x \in \bar{\xi}$
\add{\ill{}, car $\bar{\xi}$ est \ill{}} ; je dis que $\bar{\xi}$ est
$\mathrm{rel}^{\mathrm{t}}$ irréductible.
\end{enumerate}

En effet, \struck{\ill{}} \ill{} $y$ le pt.\ générique de $Y$, faisant une
extension finie $K'$ sur $K$, et normalisant $A$ dedans, d'où un $A'$ fini sur
$A$, on note que $A'$ est \emph{local} \add{et sa fibre générique
$f^{-1}(y)$ \ill{}} (car $\frac{1}{2}$ local complet $A$).

\note{le $\frac{1}{2}$ est son abréviation pour « semi- » : $A$ semi-local complet}

Utilisant le lemme
b), \struck{$\mathrm{rel}^{\mathrm{t}}$ irréductible, on \ill{}} \ill{} ; c'est
donc le cas a), \emph{absurde}.

On a donc trouvé une application $\varphi : Y \to X$ monogène : \ill{} tout
$y \in Y$ l'unique \struck{composante irréd.} $\xi \in E_{X/Y}(y)$ tel que
$x \in \bar{\xi}$, et ce $\xi$ est $\mathrm{rel}^{\mathrm{t}}$ irréd. Les
conditions (i) et (ii) pour une application monogène sont satisfaites. Il
existe donc un \ill{} $U$ de $X$ tel que $U \supset \varphi(Y)$, les fibres de
$U$ sur $Y$ \ill{} $\mathrm{rel}^{\mathrm{t}}$ irréductibles.

\page{6}

Reste à vérifier que $\varphi$ est \emph{continue}, i.e.\ que \ill{} l'image
inverse \ill{} d'un ouvert. \add{C'est vrai pour tout ouvert de $X|U$ \ill{} et
on voit \ill{}} \ill{} que $X \to \struck{Y}$ est \struck{\ill{}} ouvert.

\note{un bloc de deux lignes, enfermé dans une boucle allongée au milieu de la
page, est annulé : on y lit « $U = \varphi^{-1}(U)$ \ill{} » puis
« $M : Y \to Y$ \ill{} » ; il n'est pas rendu ici}

\textbf{Conclusion} \ill{}

\marginal{mais les composantes \struck{irréd.} $\mathrm{rel}^{\mathrm{t}}$
irréd.}

Si on ne suppose plus $f^{-1}(y)$ irréductible, \ill{} \ill{} la construction
précédente \ill{} \add{si on \ill{}} \ill{} de $X/Y$, \ill{} les composantes de
$f^{-1}(y)$, soit $Y_i = \struck{\ill{}}(Y - y)$, les diverses
$(i = 1, \ldots, n)$, et $M'$ la réunion des \ill{} ; alors les $\varphi_i :
Y'_i \to X$ \ill{} correspondent à des sections $\varphi'_i$ de $M'/Y'$, dans
\ill{} $M'$ \ill{} que les $X'_i$ \ill{} $M'$ et les $Y'_i$ \ill{}

\page{7}

On a démontré le

\textbf{Lemme.} Soit $X$ de type fini sur $Y = \mathrm{Spec}(A)$,

\begin{enumerate}
\item[(i)] $A$ \ill{} local complet ;
\item[(ii)] $f^{-1}(y)$ \struck{\ill{}} comp.\ irréductible \struck{de
$f^{-1}(y)$}, i.e.\ $\mathrm{rel}^{\mathrm{t}}$ irréductibles ;
\item[(iii)] $X$ est universellement ouvert sur $Y$ ;
\item[(iv)] \struck{la réunion} il existe $Y' = Y - y$ une
\uncertain{$Y$-préschéma} telle \struck{\ill{}} \ill{} \add{supp.\ des} \ill{}
\struck{la disparition des} composantes $\mathrm{rel}^{\mathrm{t}}$
irréductibles par $X|Y'$ ;
\item[(v)] $X \to y$ \ill{} de recollements \add{= \uncertain{coordination}} des
composantes irréductibles \ill{} = composantes de $f^{-1}(y)$, puis
\add{extension finie} de la base.
\end{enumerate}

\textbf{Alors} il \struck{existe} \struck{\ill{}} \uncertain{$Y$-préschéma} $M$
\ill{} pour le rep.\ des comp.\ $\mathrm{rel}^{\mathrm{t}}$ irréd.\ de $X/Y$.

\add{$M'/Y'$ soit séparé et fini,}

\uncertain{Le Moyen} \ill{} Supposons que \ill{} $Z$ soit \ill{} \ill{}
\struck{\ill{}} \ill{} et un \ill{} $A \to A_1$, tel que \ill{} \struck{une}
relative de $X \times_Y Y_1$, sur \ill{} \add{se spécialise en} \struck{soit}
deux distinctes \ill{}.

\textbf{Alors} $M/Y$ est séparé.

On a maintenant \struck{\ill{}} ce qu'il faut \ill{}

\note{titre souligné, un second mot biffé juste au-dessous ; la lecture
« Le Moyen » n'est pas assurée}

\page{8}

\textbf{Théorème.} $X$ \struck{\ill{} ouvert} de type fini sur $Y$. On suppose
\struck{\ill{}} que pour tout pt.\ principal $x$ d'une fibre $f^{-1}(y)$, il
existe un voisinage \ill{} \ill{} qui soit universellement ouvert et
\add{quasi-}fini sur $Y[t_1, \ldots, t_{n-1}]$. \add{Condition vérifiée p.\ ex.\
si $X$ est plat sur $Y$ au voisinage de $x$.} Pour qu'il existe un $M$
\emph{étale} et \emph{séparé} sur $Y$, \struck{\ill{}} indexant les composantes
irréductibles relatives de $X/Y$, il f.\ et s.\ que \ill{} \ill{} de relation
discrète \ill{}

\add{\ill{} les \ill{} \ill{} et \struck{\ill{}} faisant \ill{} les composantes
des relations \ill{}}, une $Y' = \mathrm{Spec}(W)$ et $X' = X \times_Y Y'$,
\ill{} fini et plat de $Y'$, \ill{} $V^{\circ}$ fini, $\xi'$ d'une \ill{} :

\begin{enumerate}
\item[(ii)] Pour toute \struck{relation} \ill{} de $f^{-1}(y')$, il existe un
pt.\ principal des \ill{} qui soit dans $\bar{\xi}'$ \add{i.e.\ une comp.\
\ill{} de la fibre générique $f^{-1}(y')$ ne se spécialise en deux composantes
$\mathrm{abs}^{\mathrm{t}}$ irréductibles distinctes de la fibre spéciale} ;
\item[(iii)] \ill{} composantes $\mathrm{abs}^{\mathrm{t}}$ irréductibles
distinctes \ill{} ne \ill{} spécialiser en \ill{} comp.\
$\mathrm{abs}^{\mathrm{t}}$ irréd.\ de $f^{-1}(y')$.
\end{enumerate}

\note{« $\mathrm{abs}^{\mathrm{t}}$ » est inséré en interligne au-dessus de
« irréductibles », deux fois ; c'est la première apparition dans le dossier de
l'abréviation d'« absolument », distincte de celle de « relativement »}

\page{9}

\textbf{N.B.} Cela signifie que \add{l'ensemble des} \struck{les composantes}
composantes irréductibles \add{relatives} de $f^{-1}(\bar{y}')$ \struck{est
\ill{}} s'injecte dans celui des composantes irréductibles
\struck{distinctes} \add{relatives} de $f^{-1}(y')$ \add{qui est \ill{}}, mais
\ill{} n'est pas néc.\ surjectif. On a cependant ceci :

\uncertain{Cor.} Pour que $M$ soit un \emph{ouvert} étale \struck{de} sur $Y$,
il f.\ et s.\ que l'on ait \ill{} les conditions précédentes :

\begin{enumerate}
\item[(iii)] \struck{Pour} Pour toute \struck{composante} \add{pt.\ générique
d'une} composante irréductible relative de $f^{-1}(y')$, il existe un pt.\
générique $x'$ d'une comp.\ irréd.\ relative de $f^{-1}(y)$ tel que
$x' \in \bar{\xi}'$.
\end{enumerate}

\page{10}

\struck{Corollaire}

\textbf{Remarques.} 1) La condition (i) \ill{} est nec.\ et suffit pour
l'existence d'un $M$ \emph{non nec.\ séparé} sur $Y$.

2) La condition (i) \struck{et (iii)} est \ill{} satisfaite si $f$ est une
application \struck{\ill{}} \add{(iii) est satisfaite. De plus,} \ill{}
\emph{séparable}. De plus, si $f$ est propre, \struck{\ill{}} \ill{} des
composantes irréductibles relatives \add{\ill{} déf. / $\kappa(y)$} \ill{} et
les conditions \ill{}. Alors \struck{la conjonction de (ii) et (iii)} \ill{} se
ramène aux conditions suivantes :

\begin{enumerate}
\item[a)] \ill{} (i) \ill{} \ill{} d'une composante $\mathrm{rel}^{\mathrm{t}}$
irréd.\ \ill{} une spécialisation \ill{} ;
\item[b)] \ill{} ;
\item[c)] $M$ existe \ill{}
\end{enumerate}

\struck{Et} Alors $M$ est propre sur $Y$.

\page{11}

\textbf{Prop.~1°} Soit $f : X \to Y$ un morphisme \struck{\ill{}}
\struck{loc.\ \ill{}} \add{de type fini} dont les fibres sont
\struck{\ill{}} irréductibles. \add{Pour tout $y \in Y$, soit $\varphi(y)$ le
pt.\ générique de $X_y$.} Conditions équivalentes :

\begin{enumerate}
\item[(i)] $\varphi$ \struck{soit} \ill{} continue ;
\item[(ii)] $f$ soit ouvert \struck{\ill{}}, \ill{} si $Y$ est loc.\ noeth.,
$f$ de type fini ;
\item[(iii)] $y \in \overline{\{y'\}} \Longrightarrow \varphi(y) \in
\overline{\varphi(\{y'\})}$.
\end{enumerate}

\textbf{Dém.} \struck{\ill{}} \ill{}

\add{(iii bis)} Pour tout $y$, $y'$ tels que \ill{} $y \in
\overline{\{y'\}}$, et $x$ au-dessus de $y$, il existe $x'$ au-dessus de $y'$
tel que $x \in \overline{\{x'\}}$.

\note{la condition (iii bis) est encadrée sur le feuillet}

En effet, \ill{} (ii) $\Rightarrow$ (iii bis) : on \ill{} $x' = \varphi(y')$,
et $x$ au-dessus de $y$ \ill{}. Il existe $x'$ \ill{} au-dessus de $y'$ tel que
\ill{}. \struck{\ill{}} (ii) $\Rightarrow$ (iii bis)), \ill{} \struck{\ill{}}
(ii)\textsuperscript{$\circ$} $\Leftrightarrow$ (ii), \ill{} alors
$\varphi^{-1}(U) = f(U)$. De plus (ii) $\Rightarrow$ (iii) banal.

\textbf{Déf.~2} On dit alors que $f$ est \ill{} \emph{universellement monogène}.
Définition de \emph{universellement monogène}.

\textbf{Déf.~3} Une \struck{\ill{}} \struck{partie} \ill{} $\varphi : Y \to X$
est dite \struck{\ill{}} \emph{monogène} s'il existe un \ill{} $U$ de $X$,
universellement monogène sur $Y$, tel que $\varphi$ soit \ill{} section de
$U/Y$.

\page{12}

\ill{} $U' \subset U$ tel que \ill{} $U' \to Y$ soit \uncertain{surjectif},
satisfaisant \struck{= propriétés} la propriété \emph{monogène}. Elle s'énonce
aussi à ceci : pour tout $y \in Y$, $\varphi(y)$ est un pt.\ générique
\add{d'une composante} de $X_y$, et il existe un \struck{\ill{}} voisinage
\ill{} tel que $U$ soit universellement ouvert \ill{} et \ill{} ses fibres
\ill{} aux pts.\ de \ill{} sont \ill{} \ill{}. Ce sont là des propriétés
locales sur $Y$.

Si $I(X/S) \to$ l'ens.\ des \ill{} monogènes de $X$ sur $S$, alors
\[
  I(X \times_S S' / S') \longrightarrow
\]
contravariant en $S'$ de façon évidente. Est-il représentable ???

\begin{tikzcd}[column sep=small, row sep=small, nodes={font=\scriptsize}]
X \arrow[dd] & & X_M \arrow[ll] \arrow[dd] & \\
& X' \arrow[ul] \arrow[dd] & & X'_{m'} \arrow[ll] \arrow[ul] \arrow[dd] \\
S & & M \arrow[ll] & \\
& S' = \mathrm{Spec}(\Omega) \arrow[ul] & & m' \arrow[ll] \arrow[ul]
\end{tikzcd}

\note{ce cube est tracé au crayon, plus pâle que le reste du feuillet ; le
montant de gauche porte deux annotations, $\kappa(s)$ en haut et $\kappa(s')$ en
bas, reliées par un trait à $S$ et à $S'$}

\page{13}

\textbf{Lemme.} Soit $f : X \to Y$ un morphisme \add{plat} de type fini
\add{en effet \ill{}}, \struck{universellement ouvert} \ill{}, avec $Y$ local
noeth. Soit \struck{$x$ le point générique} \ill{} $y \in Y$. Alors il existe un
\struck{\ill{} voisinage} $U$ \ill{} tel que $U \cap f^{-1}(y)$ soit dans
$f^{-1}(y)$ et \add{tel que pour \ill{} \struck{la réunion des composantes}}
\ill{} $y'$ \ill{} \ill{} \add{un pt.\ générique} $x'$ d'une composante de
$f^{-1}(y') \cap U$, il existe \struck{une composante} un pt.\ générique $x$
d'une composante de \struck{\ill{}} $f^{-1}(y)$ \struck{\ill{}} qui soit
spécialisation de $x'$.

\note{tout ce lemme est annulé par deux longues diagonales qui se croisent au
milieu du feuillet ; il se lit néanmoins, et il est donné}

\begin{tikzcd}[column sep=small, row sep=small]
X \arrow[d] & X' = X \times_Y Y' \arrow[l] \arrow[d] \\
Y & Y' \arrow[l]
\end{tikzcd}

$x_1 \in E_{X/Y}(y)$, \quad $x_2 \in \ill{}$, \quad
$\bar{\xi} = \mathrm{Im}(\overline{\xi'}) \ni x$

\begin{tikzcd}[column sep=small, row sep=small, nodes={font=\scriptsize}]
X \arrow[d, "f"] & X' = X \times_Y Y' \arrow[l] \arrow[d, "f'"] & X'' = X \times_Y Y'' \arrow[l] \arrow[d] \\
Y & Y' \arrow[l] & Y'' \arrow[l]
\end{tikzcd}

\note{à côté de chaque colonne il écrit les points qui l'accompagnent :
$x \leftarrow \xi$ au-dessus de $y \leftarrow \eta$ pour la première,
$x' \leftarrow \xi'$ au-dessus de $y' \leftarrow \eta'$ pour la deuxième,
$y'' \leftarrow \eta''$ pour la troisième}

\textbf{Hyp.}
\[
  \begin{cases}
    x \in E_{X/Y}(y) \\
    \xi \in E_{X/Y}(\eta) \\
    x \in \bar{\xi} \\
    y' \text{ sur } y \\
    \eta' \text{ sur } \eta \\
    y' \in \overline{\eta'}
  \end{cases}
\]

\textbf{Concl.}

\begin{enumerate}
\item[a)] Si $x' \in E_{X'/Y'}(y')$ \ill{} \add{au-dessus de} $x$, il existe
$\xi' \in E_{X'/Y'}(\eta')$ \ill{} $\xi$ tel que $x' \in \overline{\xi'}$ ;
\item[b)] Si $\xi' \in E_{X'/Y'}(\eta')$ \ill{} $\xi$, il existe
$x' \in E_{X'/Y'}(y')$ \ill{} $x$ tel que $x' \in \overline{\xi'}$.
\end{enumerate}

\note{une flèche partant de la marge gauche vient buter sur l'accolade de a) et
b) ; elle porte « Si $Y'$ \ill{} sur $Y$ et $y'$ \ill{} »}

\ill{} plus que \ill{} \ill{}. Vérifier !

\struck{Concl'}
\struck{$\Gamma_{X'/Y'}(y', \eta') \to E_{X'/Y'}(y')$}
\struck{$\Gamma_{X'/Y'}(y', \eta') \to E_{X'/Y'}(\eta')$}
\struck{\ill{} surjectif}

\note{ces deux dernières lignes, et le mot qui les suit, sont annulés par une
grande boucle}

\end{document}
